\documentclass[11pt]{article}

\usepackage[margin=1in]{geometry}
\usepackage{amsmath,amssymb,amsthm}
\usepackage{mathtools}
\usepackage[hidelinks]{hyperref}

\newtheorem{lemma}{Lemma}
\newtheorem{proposition}{Proposition}
\DeclareMathOperator{Jac}{Jac}

\begin{document}

\title{A Half of the Divisor $P_a+P_b-2\infty$ on a Genus Two Curve}
\author{}
\date{}
\maketitle

\section{Computation}This was done with GPT 5.5 Thinking.

\begin{lemma}[Quoted form of Lemma~\textup{(b)}]\label{lem:2244b}
Let
\[
X:\quad y^2=x(x+A)(x+B)(x+C)(x+D)
\]
be a genus two curve with $A=a^2$, $B=b^2$, $C=c^2$, and $D=d^2$ squares. Let $s_i$ denote the $i$-th elementary symmetric polynomial in $a,b,c,d$. Then
\[
P_0=(0,0)-\infty
\]
is divisible by $2$ in $\Jac(X)$, and one half is represented by
\[
\left[
 x^2-s_2x+s_4,
 (s_1s_2-s_3)x-s_1s_4
\right].
\]
\end{lemma}

\begin{proposition}\label{prop:half-Pa-Pb}
Let
\[
C:\quad y^2=x(x+a^2)(x+b^2)(x+c^2)(x+d^2)
\]
be a nonsingular genus two curve. Suppose that
\[
\begin{aligned}
u^2&=(a^2-c^2)(a^2-d^2),&
 v^2&=(b^2-c^2)(b^2-d^2),\\
w^2&=(a^2-c^2)(b^2-c^2),&
 t^2&=(a^2-d^2)(b^2-d^2),
\end{aligned}
\]
with $uv=wt$. Put
\[
A=a^2,\qquad B=b^2,\qquad C_0=c^2,\qquad D_0=d^2,
\]
and define
\[
p=\frac{a}{b},\qquad
r=\frac{A-C_0}{w},\qquad
s=\frac{A-D_0}{t}.
\]
Then
\[
r^2=\frac{A-C_0}{B-C_0},\qquad
s^2=\frac{A-D_0}{B-D_0},\qquad
rs=\frac{u}{v}.
\]
Let $\sigma_i$ denote the $i$-th elementary symmetric polynomial in
\[
1,\ p,\ r,\ s.
\]
Thus
\[
\begin{aligned}
\sigma_1&=1+p+r+s,\\
\sigma_2&=p+r+s+pr+ps+rs,\\
\sigma_3&=pr+ps+rs+prs,\\
\sigma_4&=prs=\frac{au}{bv}.
\end{aligned}
\]
Set
\[
\Delta=1+\sigma_2+\sigma_4=(1+p)(1+r)(1+s).
\]
Let
\[
P_a=(-a^2,0),\qquad P_b=(-b^2,0),
\]
and
\[
D=P_a+P_b-2\infty.
\]
Then $D$ is divisible by $2$ in $J(\mathbb Q)$, where $J=\Jac(C)$.
More precisely, one divisor class $D'$ satisfying
\[
2D'\sim D
\]
has Mumford representation
\[
D'=[U(x),V(x)],
\]
where
\[
U(x)=
\frac{(x+A)^2+\sigma_2(x+A)(x+B)+\sigma_4(x+B)^2}{\Delta},
\]
and $V(x)$ is the remainder modulo $U(x)$ of the cubic polynomial
\[
\frac{bv}{(A-B)^2}(x+B)^2
\left(
-\bigl(\sigma_1\sigma_2-\sigma_3\bigr)(x+A)
-\sigma_1\sigma_4(x+B)
\right).
\]
Equivalently,
\[
V(x)\equiv
\frac{bv}{(A-B)^2}(x+B)^2
\left(
-\bigl(\sigma_1\sigma_2-\sigma_3\bigr)(x+A)
-\sigma_1\sigma_4(x+B)
\right)
\pmod{U(x)},
\]
with $\deg V<2$.
\end{proposition}

\begin{proof}
Consider the change of variables
\[
X=-\frac{x+A}{x+B},\qquad
Y=\frac{(A-B)^2}{bv(x+B)^3}y.
\]
A direct calculation gives
\[
Y^2=X(X+1)(X+p^2)(X+r^2)(X+s^2).
\]
Indeed,
\[
r^2=\frac{A-C_0}{B-C_0},\qquad
s^2=\frac{A-D_0}{B-D_0},
\]
and the normalization of $Y$ uses
\[
v^2=(B-C_0)(B-D_0).
\]
Under this transformation, $P_a=(-A,0)$ maps to $(0,0)$, while
$P_b=(-B,0)$ maps to infinity. Hence the class $P_a-P_b$ on $C$ maps to
$P_0=(0,0)-\infty$ on the transformed curve. Since $2(P_b-\infty)=0$, we have
\[
P_a-P_b\sim P_a+P_b-2\infty.
\]

Applying Lemma~\ref{lem:2244b} to
\[
Y^2=X(X+1)(X+p^2)(X+r^2)(X+s^2),
\]
with square roots $1,p,r,s$, we find that one half of $P_0$ has Mumford representation
\[
[U_X(X),V_X(X)],
\]
where
\[
U_X(X)=X^2-\sigma_2X+\sigma_4
\]
and
\[
V_X(X)=(\sigma_1\sigma_2-\sigma_3)X-\sigma_1\sigma_4.
\]

We now pull this divisor back to the original $x$-coordinate. Since
\[
X=-\frac{x+A}{x+B},
\]
the equation $U_X(X)=0$ becomes
\[
0=U_X\!\left(-\frac{x+A}{x+B}\right)
=\frac{(x+A)^2+\sigma_2(x+A)(x+B)+\sigma_4(x+B)^2}{(x+B)^2}.
\]
The numerator has leading coefficient
\[
1+\sigma_2+\sigma_4=\Delta,
\]
so the normalized monic polynomial is
\[
U(x)=
\frac{(x+A)^2+\sigma_2(x+A)(x+B)+\sigma_4(x+B)^2}{\Delta}.
\]

Similarly, the relation $Y=V_X(X)$ gives
\[
\frac{(A-B)^2}{bv(x+B)^3}y
=V_X\!\left(-\frac{x+A}{x+B}\right).
\]
Therefore
\[
y=
\frac{bv(x+B)^3}{(A-B)^2}
V_X\!\left(-\frac{x+A}{x+B}\right),
\]
and hence
\[
y=
\frac{bv}{(A-B)^2}(x+B)^2
\left(
-\bigl(\sigma_1\sigma_2-\sigma_3\bigr)(x+A)
-\sigma_1\sigma_4(x+B)
\right).
\]
The Mumford representative requires the second coordinate to have degree less than
$2$, so $V(x)$ is obtained by reducing this cubic polynomial modulo $U(x)$.
Thus $[U(x),V(x)]$ is the pullback of a half of $P_0$, and consequently
\[
2[U,V]\sim P_a+P_b-2\infty.
\]
This proves the claim.
\end{proof}

\end{document}
